\documentclass[leqno]{jsarticle}
\usepackage{amssymb}
\usepackage{amsthm}
\usepackage{amsmath}
\usepackage{comment}
	%可換図式用
		%\usepackage[all]{xy}
	%重くなるので使わいときはコメント化しておく。
%定理環境 thmはあまり使わずpropをなるべく使う。
%主張は基本的に証明の中で使い、そのため番号を付けない。
%注意環境を付けてみた。
\newtheorem{thm}{定理}
\newtheorem{prop}[thm]{命題}
\newtheorem{cor}[thm]{系}
\newtheorem{lem}[thm]{補題}
\newtheorem{df}[thm]{定義}
\newtheorem{exam}[thm]{例}
\newtheorem{fact}[thm]{事実}
\newtheorem*{claim}{主張}
\newtheorem*{cau}{注意}
\renewcommand{\proofname}{{\bf 証明}}
%矢印たち。
%\toは\rightarrowと同じ。\getsは\leftarrowと同じ。同値は\iff
%上向き矢はuparrow逆はdownarrow
\newcommand{\To}{\Longrightarrow}
\newcommand{\Gets}{\Longleftarrow}
%黒板太字
\newcommand{\nn}{\mathbb{N}}
\newcommand{\zz}{\mathbb{Z}}
\newcommand{\qq}{\mathbb{Q}}
\newcommand{\rr}{\mathbb{R}}
\newcommand{\cc}{\mathbb{C}}
%無限大
\newcommand{\mgn}{\infty}
%ギリシャン
\newcommand{\dpst}{\displaystyle}
\newcommand{\ep}{\varepsilon}
\newcommand{\al}{\alpha}
\newcommand{\be}{\beta}
\newcommand{\de}{\delta}
\newcommand{\la}{\lambda}
\newcommand{\La}{\Lambda}
\newcommand{\ga}{\gamma}
\newcommand{\lil}{\la\in \La}
%商集合
\newcommand{\qt}[2]{#1/\! #2}
%enumerate環境
\renewcommand{\labelenumi}{{\rm (\arabic{enumi})}}
%以降alignじゃなくてenumerate を使え。
%なんか演算
\newcommand{\card}{\text{{\rm card}}}
\newcommand{\supp}{\text{{\rm supp}}}
%なんか演算２
\newcommand{\bd}{\text{{\rm Bdry}}}
\newcommand{\cl}{\text{{\rm CL}}}
\newcommand{\nai}{\text{{\rm Int}}}
%差集合は\setminusで統一する。等号の否定は\neq
%真部分集合は\subsetneqq　偏微分は\partial
%ディスプレイスタイル。\dfracでディスプレイ分数
%空集合は\Oを使う！！！！！！！！！！！！

\title{可算コンパクトなメタコンパクト空間はコンパクト}
\author{一色三良(concious77)}
\date{\today}
			\begin{document}
\maketitle
この文章では可算コンパクトなメタコンパクト空間がコンパクトであることを紹介する。\par
参考文献\cite{1}では$T_1$を仮定し、これを証明しているが参考文献\cite{2}では$T_1$を仮定せずとも証明できることが注意されている。\par
以下の文書では位相空間に如何なる分離公理も課さない。


%メタコンパクトの定義
\begin{df}
集合$X$の部分集合族$\mathcal{A}$と$X$の点$x$について$deg(x;\mathcal{A})$を
\[
\deg(x;\mathcal{A})=\card\{A\in \mathcal{A}\ |\ x\in A\}
\]
と定義する。これを$\mathcal{A}$に対する$x$の度数という。
\end{df}
\begin{df}
集合$X$の被覆$\mathcal{U}$について任意の点$x$について
$\deg(x;\mathcal{U})$が有限となるとき、$\mathcal{U}$を$X$の点有限被覆という。
\end{df}
\begin{df}
位相空間$X$についてその任意の開被覆に点有限被覆となる開細分が存在するとき$X$をメタコンパクト、もしくは弱パラコンパクトと言う。明らかにパラコンパクト空間はメタコンパクトである。
\end{df}


\begin{df}[既約被覆]
集合の被覆が既約被覆であるとは、その真に小さい部分族が被覆を成さないときに言う。
\end{df}


\begin{prop}
点有限被覆には既約部分被覆が存在する。
\end{prop}
\begin{proof}
集合$X$の点有限被覆を$\mathcal{U}$とし、
\[
\mathbf{C}=\{\mathcal{A}\mid \mathcal{A}\subset \mathcal{U}, \ \text{$\mathcal{A}$は$X$の被覆}\}
\]
とおく。このとき$\mathbf{C}$には包含関係により順序が備わっている。
$\mathbf{C}$の極小元が既約被覆であることに注意しよう。
以下では$\mathbf{C}$がZornの補題の仮定を充すことを確かめる。
$\mathbf{I}$を$\mathbf{C}$の全順序部分集合とする。
このとき$\bigcap\mathbf{I}$が$\mathbf{I}$の下界になることを見よう。
そのためにはこれが$X$の被覆になっていることを示せば良い。
$x\in X$を任意に取ろう。
そして$x$が属する$\mathcal{U}$の元を
$U_1,U_2, \dots, U_s\ (s=\deg(x;\mathcal{U}))$としよう。
このうちのどれか一つは$\mathbf{I}$の全てに属している。
もしそうでないとするとある$\mathcal{A}_1,\mathcal{A}_2,\dots, \mathcal{A_s}\in \mathbf{I}$が存在して、
$V_i\not\in \mathcal{A}_i$となるが、
$\mathcal{B}=\min\{\mathcal{A}_1,\mathcal{A}_2,\dots, \mathcal{A_s}\}\in \mathbf{I}$
と置くと、任意の$V_i$について$V_i\not\in \mathcal{B}$となる。
よって$x\not\in \bigcup\mathcal{B}$となり、$\mathcal{B}$が被覆であることに矛盾する。
\end{proof}

\begin{thm}[Arens-Dugundji]
可算コンパクトかつメタコンパクトな空間はコンパクトである。
\end{thm}
\begin{proof}
$X$には既約被覆$\mathcal{U}$で$\card (\mathcal{U})=\mgn$となるものが存在するとしよう。$\mathcal{U}=\{U_a\}_{a\in A}$としよう。ここで添え字付けは単射とする。$A$は無限集合なので$\{a_n\}_{n\in \nn}$がとれる。ここで添え字付けは単射とする。そして$A_n=A\setminus \{a_i\ |\ i\ge n\}$として
\[
E_n=X\setminus \bigcup_{a\in A_n}U_a
\]
と定義しよう。すると
\[
E_n\supset E_{n+1}
\]
であり、既約性から$E_n\neq\O$である。そして点有限性から
\[
\bigcap_{n\in \nn}E_n=\O
\]
なので$X$は可算コンパクトではない。以上から対偶を考えて$X$が可算コンパクトならば、既約被覆は有限被覆になる。よってメタコンパクトな可算コンパクト空間はコンパクトになる。
\end{proof}


















































	
\begin{thebibliography}{9}
\bibitem{1}R.Arens and J.Dugundji.\ \textit{Remark on the concept of compactness},\ Portugal.math.9(1950), 141-143
\bibitem{2}J.Greever,\ \textit{On Some Generalized Compactness Properties},\ Publ.RIMS. Kyoto Univ.Ser. A Vol 4(1968),39-49
\end{thebibliography}




			\end{document}



\begin{comment}
[可換図式]
\[\xymatrix{
&   & (2) \ar@{=>}[rd]&  &  &  & (5) \ar@{=>}[rd]&\\ 
& (1)\ar@{=>}[ru] \ar@{=>}[rd]&  &(4)\ar@{=>}[r]&(7)& (1)\ar@{=>}[ru] \ar@{=>}[rd]& & (7)&\\ 
&   & (3) \ar@{=>}[ur]&  &  &  &(6) \ar@{=>}[ur]&
}\]

[\bigin \endを使わない方法]

{\prop[aa] aaa}
で
\begin{prop}[aa]
aaa
\end{prop}
と同じ\df \thm \proof でも同様

[直和]
$\amalg$

[同値関係でよく使う波線]
\sim

[引数付きの新命令]
\newcommand{\prn}[1]{\|#1\|_p}

[番号のリセット]

\setcounter{equation}{0}


[字体の変更]

{\rm hogehoge}と使う。

\rm ローマン
\it イタリック
\sl 斜体

[supp]

\newcommand{\supp}{\text{{\rm supp}}}

[中央揃えの改行]

	\begin{eqnarray*}
		hoge1&=&hogehoge
		hoge2&=&hogehofe

	\end{eqnarray*}

&&の部分で揃えられる。


[\tag]
\begin{equation}
f(x) = ax^2 + bx + c \tag{1.2.3}
\end{equation}



[場合分け]

	\begin{cases}
		hoge1 & \text{状況１}\\
		hoge2 & \text{状況２}\\
		・
		・
		・
	\end{cases}



\end{comment}





